\scanpage{226}
\margin{225}
\begin{center}-- 6 --\end{center}
\margin{n$^\circ$370}
2. \uline{Le théorème d’existence et d’unicité}.

DÉFINITION 2. -- \uline{Soit} $G$ \uline{un groupe localement compact}.
\uline{On appelle mesure de Haar à gauche} (resp. \uline{à droite})
\uline{sur} $G$ \uline{une mesure positive non nulle sur} $G$,
\uline{invariante à gauche} (resp. \uline{à droite}).

THÉORÈME 1. -- \uline{Sur tout groupe localement compact},
\uline{il existe une mesure de Haar à gauche} (resp. \uline{à droite}),
\uline{et}, \uline{à un facteur constant près}, \uline{il n’en existe qu’une}.

A) \uline{Existence} -- Posons $\mathcal K(G)=\mathcal K$,
$\mathcal K_+(G)=\mathcal K_+$, $\mathcal K_+^*=\mathcal K_+-\{0\}$.
Si $C$ est une partie compacte de $G$, on notera $\mathcal K_+^*(C)$
l’ensemble des $f\in\mathcal K_+^*$ à support dans $C$.
Pour $f\in\mathcal K$ et $g\in\mathcal K_+^*$, il existe des nombres
$c_1,\ldots,c_n\ge0$ et des éléments $s_1,\ldots,s_n$ de $G$ tels que
$f\le\sum_{i=1}^n c_i\gamma(s_i)g$.
(En effet, il existe une partie ouverte non vide $U$ de $G$ telle que
$\inf_{s\in U}g(s)>0$, et le support de $f$ peut être recouvert par un nombre
fini de translatés à gauche de $U$).
Soit alors $(f:g)$ la borne inférieure des nombres $\sum_{i=1}^n c_i$
pour tous les systèmes $(c_1,\ldots,c_n,s_1,\ldots,s_n)$ de nombres $\ge0$
et d’éléments de $G$ tels que $f\le\sum_{i=1}^n c_i\gamma(s_i)g$.
On a :
\[
\begin{array}{lll}
\text{(i)} & (\gamma(s)f:g)=(f:g) & \text{pour }f\in\mathcal K,\ g\in\mathcal K_+^*,\ s\in G\ ;\\
\text{(ii)} & (\lambda f:g)=\lambda(f:g) & \text{pour }f\in\mathcal K,\ g\in\mathcal K_+^*,\ \lambda\ge0\ ;\\
\text{(iii)} & (f+f':g)\le(f:g)+(f':g) & \text{pour }f,f'\in\mathcal K,\ g\in\mathcal K_+^*\ ;\\
\text{(iv)} & (f:g)\ge(\sup f)/(\sup g) & \text{pour }f\in\mathcal K,\ g\in\mathcal K_+^*\ ;\\
\text{(v)} & (f:h)\le(f:g)(g:h) & \text{pour }f\in\mathcal K,\ g,h\in\mathcal K_+^*\ ;\\
\text{(vi)} & 0<\dfrac1{(f_0:f)}\le\dfrac{(f:g)}{(f_0:g)}\le(f:f_0)
& \text{pour }f,f_0,g\in\mathcal K_+^*\ ;
\end{array}
\]
(vii) soient $f,f',h\in\mathcal K_+$ avec $h(s)\ge1$ sur le support de $f+f'$,
et soit $\varepsilon>0$ ; il existe un voisinage compact $V$ de $e$ tel que,
pour toute $g\in\mathcal K_+^*(V)$, on ait
$(f:g)+(f':g)\le(f+f':g)+\varepsilon(h:g)$.

Les propriétés (i), (ii), (iii) sont évidentes.
Soient $f\in\mathcal K$, $g\in\mathcal K_+^*$ ;
